\section{Related work}
\label{sec:related}

\subsection{Laboratory integration and autonomous experimentation}

Laboratory integration standards seek to reduce bespoke communication and device-control effort. The original SiLA specifications defined device control, common commands, and data capture for laboratory systems \cite{bar2012sila}. SiLA~2 adopted service-oriented communication and observable properties, and associated gateway and manager projects addressed legacy devices and workflow integration \cite{juchli2022sila2,porr2020gateway,bromig2022manager}. PyLabRobot approaches interoperability through an open hardware-agnostic programming interface and community-maintained resource definitions \cite{wierenga2023pylabrobot}. These efforts establish important integration surfaces. The present study examines the remaining operational conditions that determine whether an interface can be discovered, licensed, tested, maintained, and connected to higher-level workflow state.

Autonomous laboratory literature commonly decomposes systems into experimental execution, data acquisition, model-based selection, and iterative optimization. Surveys document broad interest alongside persistent integration, standardization, and workforce constraints \cite{hung2024community}. Architectural papers describe modular autonomy and reusable orchestration \cite{canty2023autonomy,fei2024alabos}. Generalized schedulers formalize resource allocation and timing constraints \cite{cousty2024glas}. Recent deployments increasingly report maintenance, exception handling, and real-world environmental demands \cite{ochiai2025selfmaintainability,robertson2026airfree}. The forum evidence studied here complements this literature by tracing how practitioners formulate these constraints before they become formal architecture requirements or benchmark fields.

Large language models introduce another interface between generated representations and physical execution. Recent work explores direct instrumentation control through natural-language systems \cite{xie2025llmcontrol}. The forum discussions analyzed here show that command-generation competence is only one stage. Liquid behavior, sample value, third-party integration, human observation, deployment artifacts, error handling, and scientific acceptance conditions enter as additional context. The paper therefore treats AI method generation through a validation funnel that separates input handling, syntactic generation, simulation, dry execution, wet execution, and assay performance.

\subsection{Reproducibility, data stewardship, and evidence scope}

FAIR principles emphasize findability, accessibility, interoperability, and reusability for research data \cite{wilkinson2016fair}. Laboratory automation extends these concerns to executable methods, vendor databases, calibration, labware definitions, firmware, run state, and recovery records. A method file can be versioned while the active execution environment remains underidentified. The present Reproducibility Manifest Completeness metric adapts stewardship reasoning to deployment objects by asking which applicable execution fields receive immutable identifiers and which remain implicit, locally configured, or transient.

Test evidence also requires a scope discipline. Software unit tests, device simulations, dry hardware tests, wet transfer evaluations, assay validation, and production monitoring answer different questions. Forum participants explicitly contest the meaning of ``unit test'' when hardware is involved \cite{forum_unit_tests}. The paper addresses this problem through Test--Claim Alignment, which records the test object, environment, acceptance rule, observed evidence, and asserted scope. This representation allows a source-reported result such as 92 of 100 simulation traces matching source worklists to remain useful without being described as a wet-lab success rate \cite{forum_verisflow}.

\subsection{Technical communities as empirical infrastructure}

Question-and-answer communities have long supported developer knowledge exchange. Research on programming forums shows how question structure, answer practices, and social participation affect reusable technical knowledge \cite{treude2011questions,vasilescu2014social}. Open-science communities similarly reduce knowledge-access barriers by maintaining peer learning and shared practice \cite{eerland2021communities}. Specialist laboratory-automation forums combine software engineering with physical experimentation, vendor ecosystems, scientific constraints, and irreversible execution. Their discussions therefore contain evidence forms that ordinary software forums encounter less frequently, including sensor curves, liquid histories, geometry measurements, calibration, sample disposition, and wet validation.

The methodological contribution of this study lies in converting such public discourse into bounded analytical objects. The qualitative process follows explicit coding definitions and negative-case retention informed by thematic and content-analysis traditions \cite{braun2006thematic,krippendorff2018content}. Internet-research ethics motivate data minimization, contextual quotation, and careful treatment of public accessibility \cite{aoir2019ethics}. The release contains derived codes, short anonymized quotations, and provenance links. It does not redistribute a verbatim forum corpus or contributor handles.
